
\section{Burned-in Martingale Berry--Esseen}

The depth-zero martingale in \cref{lem:burn-poisson-decomp} is normalized at the
effective sample size $m = n - n_{0}$.
% We assume $m \ge n / 2$ and impose the mixing-scale burn-in lower bounds
% with logarithmic factors, so the inhomogeneous concentration term obtained
% from $\sqrt{p \, t_{\mathrm{mix}} \, \sum c_{i}^{2}}$ and then written as
% $C(u) \sqrt{p \, n}$ can be transferred to the $m$ scale.

Fix $u \in \mathbb{R}^{d}$ and put
$X_{l}^{\mathrm{bRR}} := u^{\mathsf{T}} \Delta M_{l}^{\mathrm{bRR}}$ for $2 \le l \le n - 1$. Then
\begin{equation}
|X_{l}^{\mathrm{bRR}}| \le \kappa_{\mathrm{burn}}(u),
\quad
\kappa_{\mathrm{burn}}(u)
  := 6 \, t_{\mathrm{mix}} \, C_{\mathrm{burn,Q}} \, \lVert \varepsilon \rVert_{\infty}
     \, \lVert u \rVert.
\label{eq:burn-M-incr}
\end{equation}
% This is the same bounded-increment estimate as in the stationary chapter,
% with $C_{\mathrm{burn,Q}}$ replacing the full-window weight bound.
Set
\[
s_{n,n_{0}}^{2}(u) := m \, \sigma_{n,n_{0}}^{2, \mathrm{bRR}}(u).
\]
Under \eqref{eq:burn-variance-lb-condition},
$s_{n,n_{0}}^{2}(u) \ge m \sigma^{2}(u) / 2$. Also,
$m \ge n / 2$ and the uniform bound $\lVert Q_{l}^{\mathrm{bRR}} \rVert \le C_{\mathrm{burn,Q}}$
imply the deterministic upper bound
\begin{equation}
s_{n,n_{0}}^{2}(u)
  \le 2 m \, C_{\mathrm{burn,Q}}^{2} \,
     \lVert \Sigma_{\varepsilon}^{(\mathrm{M})} \rVert \, \lVert u \rVert^{2}.
\label{eq:burn-s-upper}
\end{equation}

\begin{theorem}\label{thm:burn-M-BE}
  \textbf{(Burned-in martingale Berry--Esseen bound.)}
  Assume \textbf{UGE 1}, $\pi(\varepsilon) = 0$,
  $\lVert \varepsilon \rVert_{\infty} < \infty$, $\sigma^{2}(u) > 0$,
  $0 < \alpha$, $2 \alpha \le \alpha_{\infty}$, $m \ge n / 2$, and the
  burned-in variance lower-bound condition \eqref{eq:burn-variance-lb-condition}.
  There exist constants $C_{\mathrm{bK,1}}(u), C_{\mathrm{bK,2}}(u) > 0$, depending only on
  $\lVert u \rVert$, $\sigma(u)$, $C_{\mathrm{burn,Q}}$, $t_{\mathrm{mix}}$,
  $\lVert \varepsilon \rVert_{\infty}$, $\lVert \Sigma_{\varepsilon}^{(\mathrm{M})} \rVert$, and the
  universal constants in \cref{lem:external-martingale-be}, such that for every
  $n \ge 3$,
  \begin{equation}
d_{K} \left( \frac{u^{\mathsf{T}} M_{n,n_{0}}^{\mathrm{bRR}}}{\sqrt{m} \, \sigma_{n,n_{0}}^{\mathrm{bRR}}(u)},
    \mathcal{N}(0, 1) \right)
    \le \frac{C_{\mathrm{bK,1}}(u) \, \log^{3 / 4} n}{m^{1 / 4}}
     + \frac{C_{\mathrm{bK,2}}(u) \, \log n}{\sqrt{m}}.
\label{eq:burn-M-BE}
\end{equation}
\end{theorem}

\begin{proof} Apply \cref{lem:external-martingale-be} to the martingale differences
$X_{l}^{\mathrm{bRR}}$. The bounded-increment input is
\eqref{eq:burn-M-incr}, and the target variance is $s_{n,n_{0}}^{2}(u)$. The first and
third terms of \cref{lem:external-martingale-be} are bounded exactly as before, with
$n$ replaced by $m$ in the denominator and the harmless factor
$n / m \le 2$ absorbed into the constants. Explicitly, from
$s_{n,n_{0}}^{2}(u) \ge m \sigma^{2}(u) / 2$ and $m \ge n / 2$,
\begin{equation}
\frac{(2 n + 1) \log\left(2 n + 1\right)}{s_{n,n_{0}}^{3}(u)}
  \le C(u) \frac{\log n}{\sqrt{m}}.
\label{eq:burn-martingale-be-first-nm}
\end{equation}
The bounded-increment Lindeberg term is controlled in the same way, using
\eqref{eq:burn-s-upper} and $s_{n,n_{0}}^{2}(u) \ge m \sigma^{2}(u) / 2$. For the
conditional-variance term use
\cref{lem:burn-bracket-conc}:
\[
\mathbb{E}^{1 / p} \left[ | u^{\mathsf{T}} \langle M^{\mathrm{bRR}} \rangle_{n,n_{0}} u
    - s_{n,n_{0}}^{2}(u) |^{p} \right]
  \le C \, \sqrt{p \, n}.
\]
% Here the constant contains the explicit bracket coefficient
% $C_{\mathrm{burn,Q}}^{2} \lVert u \rVert^{2} \lVert \varepsilon \rVert_{\infty}^{2} t_{\mathrm{mix}}^{5 / 2}$
% from \cref{lem:burn-bracket-conc}; the underlying concentration input is
% \eqref{eq:external-markov-conc}.
Together with $m \ge n / 2$, \eqref{eq:burn-variance-lb-condition}, and
\eqref{eq:burn-s-upper}, this gives
\[
\mathrm{Term II}
  \le C(u) \, p^{(3 p + 1) / (2 (2 p + 1))}
     \, m^{-p / (2 (2 p + 1))}.
\]
Taking $p = \left\lceil \log n \right\rceil$ gives the displayed
$\log^{3 / 4}(n) m^{-1 / 4}$ bound. The bounded-increment terms give
the second displayed term. \end{proof}
